\section{Related Work}
\label{sec:related-work}

\subsection{Dataset Distillation}
Dataset distillation synthesizes a compact dataset $\mathcal{S}$ such that a model trained on $\mathcal{S}$ approximates performance on the full dataset $\mathcal{T}$. The field was formalized through gradient matching~\cite{zhao2023improved}, trajectory matching~\cite{cazenavette2022dataset}, and distribution matching~\cite{liu2023dream}, which jointly optimize synthetic data and a surrogate model in a bi-level loop. While effective on small benchmarks (CIFAR, SVHN), these methods scale poorly to ImageNet~\cite{russakovsky2014imagenet} due to the prohibitive bi-level cost. SRe$^2$L~\cite{yin2023squeeze} broke this barrier by decoupling optimization: train a generative model, synthesize unlimited images, then train a classifier. IDC~\cite{kim2022dataset} improved synthetic-data parameterization through differentiable augmentation, and the DD-Ranking benchmark~\cite{li2025ddranking} provided standardized evaluation, revealing that many methods struggle to outperform simple real-data subsampling.

\subsection{Generative Dataset Distillation with Diffusion Models}
Pretrained diffusion models~\cite{ho2020denoising, peebles2022scalable} have opened a new frontier. D$^4$M~\cite{su2024dataset} disentangled label-conditional and image-conditional generation using Stable Diffusion. MGD$^3$~\cite{chansantiago2025modeguided} directly samples from a frozen DiT-XL/2 on ImageNet-1K, achieving state-of-the-art performance without any fine-tuning. The field has expanded rapidly: CoDA~\cite{zhou2025coda} exploited text-to-image models for training-free distillation, PRISM~\cite{moser2025prism} decoupled architectural priors, and LGD~\cite{chansantiago2026learnabilityguided} introduced curriculum-based sample selection. Most recently, DMGD~\cite{wang2026dmgd} proposed distribution matching, SAS~\cite{li2026semanticaware} introduced semantic-aware latent selection, ManifoldGD~\cite{roy2026manifoldgd} applied hierarchical manifold guidance, and Pool-Select-Refine~\cite{li2026poolselectrefine} introduced allocation-aware distillation. Our work builds on MGD$^3$'s training-free paradigm but replaces its fixed guidance with per-class adaptive scheduling, occupying a complementary position to LGD's training-based approach.

\subsection{Classifier-Free Guidance and Adaptive Guidance Schedules}
CFG~\cite{ho2022classifierfree} has become the de facto conditioning mechanism for diffusion models, enabling post-hoc control of the guidance strength $\lambda$ at inference time. Recent work has recognized that optimal guidance varies across timesteps and conditioning inputs: \citet{yehezkel2025navigating} introduced annealed guidance scales, \citet{zhang2025much} showed guidance should be modulated by conditioning complexity, and \citet{malarz2025classifierfree} proposed adaptive scaling based on timestep and conditioning. However, in the dataset distillation context, these insights remain unexplored: MGD$^3$ uses a single fixed $\lambda$, and no prior work has systematically studied how per-class guidance adaptation affects downstream classification. \ags{} bridges this gap by introducing a per-class complexity metric specifically designed for the distillation objective, where the target is discriminative feature preservation rather than aesthetic fidelity.
